\documentclass[11pt]{article}

\usepackage[a4paper,margin=0.82in]{geometry}
\usepackage{amsmath,amssymb,amsthm,mathtools}
\usepackage{microtype}
\usepackage[hidelinks]{hyperref}

\newtheorem{theorem}{Theorem}
\newtheorem{lemma}{Lemma}
\theoremstyle{remark}
\newtheorem{remark}{Remark}

\newcommand{\Z}{\mathbb Z}
\newcommand{\N}{\mathbb N}

\title{An Infinite Pell Family for
\texorpdfstring{$y^2+x^2y+xz^2-3=0$}{y2 + x2 y + x z2 - 3 = 0}}
\author{}
\date{}

\begin{document}
\maketitle
\vspace{-2.2em}

\begin{abstract}
We give a completely explicit infinite family of integer solutions to
\[
 y^2+x^2y+xz^2-3=0.
\]
The construction is elementary: fixing $x=-2$ reduces the equation to the generalized Pell equation
$(y+2)^2-2z^2=7$.  A single solution of the latter is propagated by an explicit norm-preserving recurrence.  No general theorem about Pell equations is required.
\end{abstract}

\section{The infinite family}

\begin{lemma}[Norm-preserving recurrence]\label{lem:recurrence}
Let $(U_n,V_n)_{n\geq 0}$ be defined by
\[
 (U_0,V_0)=(3,1),
 \qquad
 \begin{cases}
 U_{n+1}=3U_n+4V_n,\\
 V_{n+1}=2U_n+3V_n.
 \end{cases}
\]
Then, for every $n\geq 0$,
\[
 U_n^2-2V_n^2=7.
\]
Moreover, $U_n$ and $V_n$ are positive, and $(V_n)_{n\geq0}$ is strictly increasing.
\end{lemma}

\begin{proof}
For arbitrary integers $U,V$ one has the identity
\begin{align*}
 (3U+4V)^2-2(2U+3V)^2
 &=9U^2+24UV+16V^2 \\
 &\quad{}-2(4U^2+12UV+9V^2)\\
 &=U^2-2V^2.
\end{align*}
Thus the recurrence preserves the quadratic form $U^2-2V^2$.  Since
$U_0^2-2V_0^2=9-2=7$, induction gives $U_n^2-2V_n^2=7$ for all $n$.

The recurrence has integral coefficients and integral initial values, so $U_n,V_n\in\Z$ for all $n$.  The initial values are positive, and the recurrence has positive coefficients, so $U_n,V_n>0$ for all $n$.  Finally,
\[
 V_{n+1}-V_n=2U_n+2V_n>0,
\]
so $V_{n+1}>V_n$ for every $n$.
\end{proof}

\begin{theorem}\label{thm:main}
For every integer $n\geq0$, the two triples
\[
 \boxed{(x,y,z)=\bigl(-2,\,U_n-2,\,\pm V_n\bigr)}
\]
are integer solutions of
\[
 y^2+x^2y+xz^2-3=0.
\]
Consequently, the equation has infinitely many integer solutions.
\end{theorem}

\begin{proof}
Fix $n\geq0$ and put $x=-2$, $y=U_n-2$, and $z=\pm V_n$.  By Lemma~\ref{lem:recurrence},
\begin{align*}
 y^2+x^2y+xz^2-3
 &=(U_n-2)^2+4(U_n-2)-2V_n^2-3\\
 &=U_n^2-2V_n^2-7\\
 &=0.
\end{align*}
Hence every displayed triple is a solution.  The positive integers $V_n$ are strictly increasing by Lemma~\ref{lem:recurrence}; therefore the triples with $z=V_n$ are pairwise distinct.  This proves infinitude.
\end{proof}

The first values are
\[
 (U_n,V_n)=(3,1),(13,9),(75,53),(437,309),\ldots,
\]
which yield
\[
 (-2,1,\pm1),\quad
 (-2,11,\pm9),\quad
 (-2,73,\pm53),\quad
 (-2,435,\pm309),\ldots.
\]

\section{How the construction was found}\label{sec:discovery}

The original expression is quadratic in $y$.  Completing the square is therefore the natural first move:
\begin{align}
 4\bigl(y^2+x^2y+xz^2-3\bigr)=0
 \quad\Longleftrightarrow\quad
 (2y+x^2)^2+4xz^2=x^4+12.\label{eq:square}
\end{align}
The sign of $x$ now becomes decisive.  When $x>0$, the two square terms on the left of \eqref{eq:square} have the same sign, so this specialization does not naturally produce an unbounded indefinite quadratic equation.  To obtain a Pell-type mechanism, one instead takes $x<0$.  Writing $x=-m$ with $m>0$ gives
\begin{equation}
 (2y+m^2)^2-4mz^2=m^4+12.\label{eq:m}
\end{equation}

The smallest choices illustrate what is needed.  For $m=1$, equation \eqref{eq:m} becomes
\[
 (2y+1)^2-(2z)^2=13,
\]
a difference of two integral squares.  It factors as
$(2y+1-2z)(2y+1+2z)=13$, so this route can give only finitely many pairs.

The next choice, $m=2$, is qualitatively different.  Equation \eqref{eq:m} becomes
\[
 (2y+4)^2-8z^2=28,
\]
and division by $4$ gives
\begin{equation}
 (y+2)^2-2z^2=7.\label{eq:pell}
\end{equation}
This is an indefinite quadratic equation with the obvious seed
\[
 3^2-2\cdot1^2=7.
\]
What remains is to reproduce this seed indefinitely without changing the value $7$.

For this purpose, write a pair $(U,V)$ formally as $U+V\sqrt2$ and define its norm by
\[
 N(U+V\sqrt2)=U^2-2V^2.
\]
Indeed, since
\[
 (U+V\sqrt2)(A+B\sqrt2)=(UA+2VB)+(UB+VA)\sqrt2,
\]
one obtains by direct expansion
\begin{align*}
 N\bigl((U+V\sqrt2)(A+B\sqrt2)\bigr)
 &=(UA+2VB)^2-2(UB+VA)^2\\
 &=(U^2-2V^2)(A^2-2B^2)\\
 &=N(U+V\sqrt2)N(A+B\sqrt2).
\end{align*}  Hence multiplication by any element of norm $1$ preserves equation \eqref{eq:pell}.  The small identity
\[
 3^2-2\cdot2^2=1
\]
provides such an element, namely $3+2\sqrt2$.  Multiplying gives
\[
 (U+V\sqrt2)(3+2\sqrt2)
   =(3U+4V)+(2U+3V)\sqrt2,
\]
which is exactly the recurrence in Lemma~\ref{lem:recurrence}.  Starting from $3+\sqrt2$, whose norm is $7$, repeated multiplication by $3+2\sqrt2$ produces
\[
 U_n+V_n\sqrt2=(3+\sqrt2)(3+2\sqrt2)^n.
\]
The recurrence proof in Section~1 then establishes integrality, preservation of the norm, and strict growth without appealing to any external result on Pell equations.

\end{document}
